\section{Discussion}\label{sec:discussion}

\subsection{Why Dense Physical Graphs Fail}

The most striking finding is that removing the physical graph entirely improves RMSE by 32.8\% (Table~\ref{tab:oversmoothing}). This oversmoothing effect arises because the physical graph's high density (2833 edges for 207 nodes) causes repeated graph convolution operations to aggregate information from increasingly distant and irrelevant nodes, washing out the distinctive features of individual road segments. This finding is consistent across both datasets and aligns with theoretical predictions about over-squashing in graph neural networks.

\subsection{Why Multi-Lag Processing Helps}

The multi-lag formulation provides two distinct advantages. First, it makes the temporal dependency structure explicit: instead of learning a single ambiguous graph, we obtain separate graphs for each temporal lag, each with a clear interpretation (Table~\ref{tab:lag_stats}). Second, the T-GCN-MultiGSL-Mix architecture can process these graphs through separate T-GCN cells and combine the results through learned mixing, preserving the lag-specific information that would be lost by merging into a single adjacency.

The fact that T-GCN-MultiGSL (same edges, fixed assignment) achieves 4.794 while T-GCN-MultiGSL-Mix achieves 4.452 demonstrates that \textit{how} the lag-specific graphs are processed matters as much as \textit{which} graphs are used. The 0.342 RMSE improvement from learned graph mixing represents a 7.1\% further reduction beyond the fixed multi-graph approach. Training converges stably within 50 epochs for all methods, with T-GCN-MultiGSL-Mix reaching the lowest final training loss (Appendix~\ref{app:diagnostics}).

\subsection{Why the Effect is Dataset-Dependent}

The improvement on SZ-Taxi is marginal ($\leq0.3\%$ at PH=1--3 and $-0.02\%$ at PH=4, single seed) compared to the strong improvement on Los-loop (14.9\% five-seed mean). Two hypotheses can be ruled out. First, temporal resolution is not the explanation: SZ-Taxi and the 15-minute Los-loop variant share the same sampling interval, yet Los-loop improves strongly at both 5-minute and 15-minute sampling (Sections~\ref{sec:multiseed} and~\ref{sec:los15min}), so the divergence tracks the dataset rather than the sampling interval. Second, capacity is not the explanation (Section~\ref{sec:param_control}). A candidate mechanism consistent with the evidence is graph learnability: at the operating threshold, DAGMA retains only 2 cross-sensor edges for SZ-Taxi versus 30--32 for Los-loop, suggesting that the urban network's dependency structure offers the learner far less exploitable signal than highway propagation patterns, which have clear temporal structure (vehicles traveling at highway speeds create predictable lagged dependencies). Other dataset differences (156 vs.\ 207 sensors, urban vs.\ highway dynamics) may also contribute. This dataset dependence is an honest finding that the paper explicitly reports rather than hiding.

\subsection{Connection to Original GSL/cGSL Findings}

The original GSL/cGSL experiments established that learned sparse graphs outperform dense physical graphs, and that different backbone architectures (GCN vs T-GCN) benefit from different graph structures. These findings remain valid and motivated the deeper investigation into temporal heterogeneity. The multi-lag formulation can be viewed as a principled extension: rather than debating whether the learned graph should be cyclic or acyclic, we construct lag-specific graphs that naturally capture different temporal aspects of the dependency structure.


